\section{詳細解説}

\subsection{スイング・ツイスト分解}

\begin{figure}[htbp]\centering
\resizebox{0.9\linewidth}{!}{% スイング・ツイスト分解
\begin{tikzpicture}[>=Stealth,font=\small]
\begin{scope}
  \fill (0,0) circle (2pt) node[left]{肩};
  \draw[line width=4pt,gray!50] (0,0)--(0,-2);
  \draw[line width=4pt,blue!60] (0,0)--(1.4,-1.4);
  \draw[->,thick] (0,-1.2) arc (-90:-45:1.2);
  \node at (0.75,-1.35) {スイング};
  \node[align=center,font=\footnotesize] at (0.6,-2.6) {腕の向きを変える\\（2 自由度）};
\end{scope}
\begin{scope}[shift={(4.5,0)}]
  \fill (0,0) circle (2pt);
  \draw[line width=4pt,blue!60] (0,0)--(1.4,-1.4);
  \draw[thick] (0.9,-0.9) ellipse (0.35 and 0.15);
  \draw[->,thick] (1.25,-0.9) arc (0:-200:0.35 and 0.15);
  \node at (1.8,-0.6) {ツイスト};
  \node[align=center,font=\footnotesize] at (0.6,-2.6) {腕の長軸まわりにひねる\\（1 自由度）};
\end{scope}
\node[font=\footnotesize,align=left] at (9.2,-1.0) {上腕の姿勢 $R = R_{\text{swing}}R_{\text{twist}}$\\ひねりの値は分解の仕方\\（基準姿勢・順序）に依存する};
\end{tikzpicture}
}
\caption{スイング・ツイスト分解：上腕の姿勢を、腕の向きを変える回転とひねりに分ける。}\label{fig:tikz-swingtwist}
\end{figure}

上腕の姿勢（回転 $R$）を、「腕の向きを変える回転（スイング）」と「腕の長軸まわりのひねり（ツイスト）」に分けて表すことを、スイング・ツイスト分解という。
\begin{equation}
R = R_{\text{swing}}\,R_{\text{twist}},\qquad R_{\text{twist}}=\text{（上腕の基準姿勢での長軸まわりの回転）}.
\end{equation}
四元数 $(x,y,z,w)$ で表した $R$ から、ツイスト角は長軸方向の成分 $v_a$ を使って $2\,\mathrm{atan2}(v_a, w)$ で求まる。
ゲームの物理エンジンで、人型キャラクターの肩の可動域を「円錐（スイングの範囲）＋ひねりの範囲」で表すのと同じ考え方である。

ただし、\textbf{ひねりの値は分解の仕方（基準姿勢、回転の順序）に依存する}。コッドマンのパラドックスはその典型で、
同じ最終姿勢でも、どの経路でそこに来たかによって「ひねった量」の解釈が変わる。
バイオメカニクスで肩の角度の測り方に国際的な標準 \cite{R36} が定められているのは、このためである。

\subsection{立体角の例題}

\begin{example}
腕を真下から前方 90°（水平）に上げ、そのまま水平面内で外側（横）へ 90° 回し、真下に下ろす。
腕の向きが描く球面三角形は、3 つの頂点（下・前・横）がそれぞれ互いに 90° 離れた「八分球」で、立体角は全球 $4\pi$ の 1/8、つまり $\pi/2$ である。
したがって上腕は 90° ひねられる。腕を半頂角 45° の円錐を描くように回して下ろすと、立体角は $2\pi(1-\cos45^\circ) = 1.840$ rad（105.4°）になる。
\end{example}

\subsection{胸の厚さと必要な内旋}

胸の奥行を 0.16〜0.28 m で細かく変えて、腕組みに必要な内旋の最小値を求めた（図~\ref{fig:D11_chest_sweep.png}）。
\begin{longtable}{rrrr}
\toprule
胸の奥行 & 球関節のみ & 球関節＋上肢帯 & 使った前方突出 \\ \midrule \endhead
0.16 m & 83° & 70° & 11° \\
0.20 m & 89° & 72° & 14° \\
0.22 m & 93° & 73° & 16° \\
0.24 m & 97° & 74° & 18° \\
0.26 m & 110° & 75° & 20° \\
0.28 m & \textbf{解なし} & 77° & 22° \\
\bottomrule
\end{longtable}
\begin{itemize}
\item 球関節のみでは、胸の奥行 0.20 m を超えると人の範囲（90°）を超え、0.28 m では腕組みそのものができなくなった（水平内転の上限 45° の下で、どのひねりでも届かない）。
\item 上肢帯の前方突出を使うと、必要な内旋はどの厚さでも 70〜77° に収まる。その代わり、使う前方突出は胸が厚いほど大きくなる。
\end{itemize}
\textbf{胸の厚さ（装甲・筋肉）と肩の可動域はトレードオフの関係にあり、厚い胸には肩甲骨の機構が必須になる}。

\DFIG{D11_chest_sweep.png}{胸の奥行と、腕組みに必要な内旋（左軸）・使った前方突出（右軸、緑）。灰色は人の内旋の可動域の目安。}

\subsection{腕組みの判定の全結果}

\DTAB{sim08_arm_crossing}

\subsection{機械設計への示唆：関節の種類}

\begin{longtable}{p{0.22\linewidth}p{0.34\linewidth}p{0.34\linewidth}}
\toprule
関節の種類 & 長所 & 腕組みへの影響 \\ \midrule \endhead
2 軸ジンバル（旋回＋俯仰） & 構造が簡単、関節角が直観的 & ひねりを選べず不可能。極で特異点 \\
3 軸（球関節に相当） & ひねりを選べる & 胸が薄ければ可能だが、人の限界ぎりぎりの内旋が必要 \\
3 軸＋上肢帯（肩の付け根を動かす） & 人に近い可動域 & 胸が厚くても可能。関節が増え、制御が複雑になる \\
\bottomrule
\end{longtable}

\section*{考察課題}
\begin{exercise}
腕組みのほかに、人間らしい姿勢で肩甲骨の動きが必要になるものを挙げ（例：背中で手を組む、頭の後ろで手を組む）、それぞれ必要な肩の動きを考えよ。
\end{exercise}
\begin{exercise}
モーションキャプチャーで計測した人の腕の動きを 3 軸の肩のロボットに移すとき、コッドマンのパラドックスによってどのような誤差が生じうるか。対策を考えよ。
\end{exercise}
